
\begin{thebibliography}{*****}

\bibitem[Ab]{Ab} W. Abikoff. The real analytic theory of Teichm\"uller space. 
      Lect. Notes Math, v. 820, Springer, Berlin. 

\bibitem[A1]{A} L. Ahlfors. Conformal invariants: Topics in geometric function theory. McGraw Hill, 1973.

\bibitem[A2]{Ahlfors:quasiconformal}
L. Ahlfors.  Lectures on quasiconformal mappings.
 The Wadsworth \& Brooks/Cole Mathematics Series. Wadsworth \&
  Brooks/Cole Advanced Books \& Software, Monterey, CA, 1987.
  
\bibitem[A3]{Ahlfors:complex}
L. V. Ahlfors. Complex analysis, third edition. McGraw-Hill, 1979.

\bibitem[AL]{AL} A. Avila \& M. Lyubich.
Hausdorff dimension and conformal measures of Feigenbaum Julia sets. Journal of the American Mathematical Society. v. 21, no. 2 (2008), pp. 305-63.



\bibitem[DH1]{Orsay} A. Douady \& J.H. Hubbard.
  \'{E}tude dynamique des polyn\^{o}mes complexes.
 Publication Mathematiques d'Orsay,
84-02 and 85-04.


\bibitem[DH2]{DH}
 A. Douady \& J.H. Hubbard. On the dynamics of polynomial-like maps.
    Ann. Sc. \'{E}c. Norm. Sup., v. 18 (1985), 287-343.

\bibitem[DH3]{DH-Thurston}  A. Douady \& J.H. Hubbard.
  A proof of Thurston's topological characterization of rational functions.
  Acta Math., v. 171 (1993), 263--297. 
  
\bibitem[DL]{satellite} D. Dudko \& M. Lyubich.
MLC at Feigenbaum points. arXiv:2309.02107


\bibitem[F]{Fenchel:Elementary}
W. Fenchel.
 Elementary geometry in hyperbolic space.
  Gruyter Studies in Mathematics, v. 11.
   Walter de Gruyter \& Co., Berlin, 1989.

\bibitem[G]{Gamelin}
Gamelin, Theodore. Complex analysis. Springer Science \& Business Media, 2003.

\bibitem[H]{H} J.H. Hubbard.  Local connectivity of Julia sets and bifurcation loci:
      three theorems of J.-C. Yoccoz. 
     In: ``Topological Methods in Modern Mathematics, A Symposium in
     Honor of John Milnor's 60th Birthday", Publish or Perish, 1993.

\bibitem[HJ]{HJ} J. Hu \& Y. Jiang. The Julia set of the Feigenbaum quadratic polynomial is
  locally connected. Preprint 1993. 

\bibitem[J]{J} Y. Jiang. Infinitely renormalizable quadratic polynomials. 
    Trans. AMS, v. 352 (2000), 5077-5091. 






\bibitem[KL1]{covering lemma} J. Kahn \& M. Lyubich.
The quasi-additivity law in conformal geometry. Annals of mathematics (2009), pp.561-593.

\bibitem[KL2]{decorations} J. Kahn \& M. Lyubich. {\it A priori bounds} for some infinitely
renormalizable quadratics: II.  Decorations.
Ann. Scient. \'Ec. Norm. Sup., v.~41 (2008), 57 -- 84.

\bibitem[KL3]{molecules} J. Kahn \& M. Lyubich. {\it A priori bounds} for some infinitely
renormalizable quadratics: III.  Molecules.  In: ``Complex Dynamics: Families and Friends''.
    Proceeding of the conference dedicated to Hubbard's 60th birthday (ed.:  D. Schleicher). 
    Peters, A K, Limited, 2009. 

\bibitem[LS]{LS} G. Levin \& S. van Strien.
    Local connectivity of Julia sets  of real polynomials.
   Ann. Math., v. 147 (1997), 471-541.  


\bibitem[L1]{puzzle} M. Lyubich. Dynamics of quadratic polynomials, I-II. 
   Acta Math., v. 178 (1997), 185 -- 297.

\bibitem[L2]{universe} M. Lyubich. 
  Feigenbaum-Collet-Tresser Universality and Milnor's Hairiness Conjecture.
  Ann. Math., v. 156 (1999), 1-78. 

\bibitem[L3]{encycl} M. Lyubich. Renormalization and Universality. 
In: Encyclopedia of Math. Sci., Elsevier 2006. 

\bibitem[LY]{LY} M. Lyubich \& M. Yampolsky. 
    Complex bounds for real maps. Ann. Inst. Fourier, v. 47 (1997), 1219-1255. 







\bibitem[M]{M} J. Milnor. Self-similarity and hairiness in the Mandelbrot set.
 In: Computers in Geometry and Topology, 
 Lecture Notes in Pure and Applied Math., v. 114 (1989), 211-257. 

\bibitem[McM1]{McM} C. McMullen. Complex dynamics and renormalization.
   Annals of Math. Studies, v. 135.
   Princeton University Press, 1994. 

\bibitem[McM2]{towers} C. McMullen.
   Renormalization and 3-manifolds 
which fiber over the circle.  Annals of Math. Studies, v. 142.
 Princeton University Press, 1996. 


\bibitem[MS]{MS} W. de Melo \& S. van Strien. 
   One dimensional dynamics. Springer-Verlag, 1993.

\bibitem[S]{S} D. Sullivan.
     Bounds, quadratic differentials,
and renormalization conjectures.  AMS Centennial Publications,
v. 2: Mathematics into Twenty-first Century (1992).  


\bibitem[P]{P} K. Pilgrim. Thesis. 

\end{thebibliography}




